\section{Computational Results}
\label{sec:results}

We implemented the Spin Gem framework in Python and JavaScript, conducting experiments to validate the theoretical predictions.

\subsection{The Spin Gem Concept}

Figure~\ref{fig:concept} illustrates the Spin Gem construction for a hexagonal planar molecule. The 3D visualization shows atoms colored by spin value with arrows indicating spin direction. The covariance matrix reveals the spin-position coupling structure.

\begin{figure}[H]
    \centering
    \includegraphics[width=\textwidth]{figures/fig1_spin_gem_concept.png}
    \caption{\textbf{Spin Gem construction.} (a) Hexagonal molecule with alternating spins. (b) Point cloud in $(x, y, s)$ space. (c) The $4 \times 4$ covariance matrix. This visualization is available interactively in Tab 1 of the web demo.}
    \label{fig:concept}
\end{figure}

\subsection{Magic vs.\ Random Configurations}

Figure~\ref{fig:magic_random} compares energy distributions for random spin assignments on symmetric geometries. For each geometry, we sampled 20,000 random spin configurations with $s_k \in [-1, 1]$.

\begin{figure}[H]
    \centering
    \includegraphics[width=\textwidth]{figures/fig2_magic_vs_random.png}
    \caption{\textbf{Energy distributions for symmetric geometries.} Top: Probability density (log scale). Bottom: Minimum-energy configurations. Magic configurations ($E=0$) are rare in all cases. The web demo (Tab 3) allows interactive exploration with live histograms.}
    \label{fig:magic_random}
\end{figure}

\subsection{The Magic Square Connection}

Figure~\ref{fig:magic_square} demonstrates the bridge between magic squares and Spin Gems. The Lo Shu square achieves exactly $\text{Cov}(X, S) = \text{Cov}(Y, S) = 0$.

\begin{figure}[H]
    \centering
    \includegraphics[width=\textwidth]{figures/fig3_magic_square_connection.png}
    \caption{\textbf{Magic squares as Spin Gems.} (a) Lo Shu square. (b) Corresponding Spin Gem. (c) Covariance scatter plot: 500 random arrangements (gray) vs.\ Lo Shu (red star at origin). See Tab 2 of the web demo for interactive exploration.}
    \label{fig:magic_square}
\end{figure}

\subsection{Precession Dynamics}
\label{sec:dynamics}

Figure~\ref{fig:precession} shows spin precession in an external magnetic field $\mathbf{B} = B_0 \hat{\mathbf{z}}$.

\begin{figure}[H]
    \centering
    \includegraphics[width=\textwidth]{figures/fig4_precession_dynamics.png}
    \caption{\textbf{Precession dynamics.} Top: Magic configuration. Bottom: Random configuration. Magic configurations exhibit more stable energy evolution during precession. Tab 5 of the web demo provides real-time animation.}
    \label{fig:precession}
\end{figure}

\subsection{Energy Landscape Structure}

Figure~\ref{fig:landscape} visualizes the energy surface for the tetrahedron, parameterized by two spins.

\begin{figure}[H]
    \centering
    \includegraphics[width=\textwidth]{figures/fig5_energy_landscape_3d.png}
    \caption{\textbf{Energy landscape.} (a) Surface plot. (b) Contour plot with magic point marked. (c) Radial energy profile. The 3D surface is rendered interactively in Tab 4 of the web demo.}
    \label{fig:landscape}
\end{figure}

\subsection{Scaling with Atom Number}

Figure~\ref{fig:scaling} examines how energy scales with the number of atoms.

\begin{figure}[H]
    \centering
    \includegraphics[width=\textwidth]{figures/fig6_scaling_atoms.png}
    \caption{\textbf{Scaling behavior.} (a) Mean and minimum energy vs.\ $N$. (b) Energy distributions for $N = 4, 8, 12, 20$. (c) Probability of near-magic configurations decreases with $N$.}
    \label{fig:scaling}
\end{figure}

\subsection{Summary}

\begin{table}[H]
\centering
\begin{tabular}{lcccc}
\toprule
Geometry & $N$ & Mean $E$ & Min $E$ (sampled) & Magic Exists? \\
\midrule
Tetrahedron & 4 & 0.082 & $2.1 \times 10^{-4}$ & Yes \\
Octahedron & 6 & 0.051 & $8.7 \times 10^{-5}$ & Yes \\
Cube & 8 & 0.043 & $3.2 \times 10^{-5}$ & Yes \\
$3 \times 3$ lattice & 9 & 0.21 & 0 (exact) & Yes (magic square) \\
\bottomrule
\end{tabular}
\caption{Summary statistics for spin balance energy.}
\label{tab:summary}
\end{table}
